\section{Background of the Study}
Communication is a fundamental requirement for social participation, education, healthcare access, and economic opportunity. For members of the deaf and hard-of-hearing community in India, Indian Sign Language (ISL) is a natural and expressive medium that encodes meaning through coordinated hand shape, hand orientation, movement trajectory, body posture, and facial expression. In practice, however, smooth communication between ISL users and non-signers is still limited by the availability of trained interpreters and by the absence of reliable real-time assistive tools in many public contexts.

The research community has shown that deep learning methods can significantly improve visual recognition tasks, including gesture and sign understanding. Modern computer vision models can learn robust spatial representations from raw RGB frames, while sequence models can capture temporal evolution across consecutive frames. These developments are strongly aligned with sign language understanding, where meaning depends not only on static hand configuration but also on motion continuity and contextual timing \cite{lecun2015deep,camgoz2018neural,joze2019msasl}.

In the Indian deployment context, the challenge is not only recognition quality but also operational stability. A practical system must sustain consistent behavior under variable lighting, camera motion, and signer-specific style differences. It must also remain responsive enough for conversational use, where delay directly affects communication flow and user trust.

Therefore, the background of this study is framed around a dual objective: achieving robust visual-temporal understanding and ensuring real-time usability. This framing motivates an architecture that combines efficient feature extraction, temporal modeling, and selective attention so the system can remain accurate without becoming too heavy for practical inference.

Despite this progress, ISL-focused systems still face practical and technical constraints. Gesture classes can be visually similar, signer-specific variations can be large, and unconstrained environments introduce noise from lighting, background clutter, and camera perspective. Therefore, there is a clear need for an ISL recognition framework that is both accurate and computationally efficient, and that can operate with real-time responsiveness for deployment-oriented scenarios.
